\documentclass[11pt]{article}

\usepackage[utf8]{inputenc}
\usepackage[T1]{fontenc}
\usepackage{amsmath,amssymb,amsfonts}
\usepackage{booktabs}
\usepackage{array}
\usepackage{geometry}
\usepackage{hyperref}
\usepackage{enumitem}
\usepackage{xcolor}
\usepackage{microtype}
\geometry{margin=1in}
\hypersetup{colorlinks=true,linkcolor=blue,citecolor=blue,urlcolor=blue}

\newcommand{\dimname}[1]{\textsc{#1}}
\newcommand{\passk}{$\mathrm{pass}^{k}$}

\title{LAIBench: A Governance-Oriented Benchmark for\\
Critical-Finding Safety and Reliability in\\
AI Radiology Report Generation}

\author{LAIBench Authors\thanks{Author list, affiliations, and corresponding
contact are pending confirmation before any formal submission. This draft is
released for technical review.}}
\date{Draft version --- May 2026}

\begin{document}
\maketitle

\begin{abstract}
Mainstream metrics for AI radiology report generation optimize for average
linguistic or entity-overlap quality and report a single aggregate number. That
framing is poorly matched to the governance question an institution actually
faces before deployment: \emph{does this system reliably preserve the
information it was handed, every time, without dropping or inverting a critical
finding?} We present LAIBench, a benchmark framework that treats radiology
report generation as a faithfulness-and-reliability problem rather than a
quality-maximization problem. We are explicit about scope: LAIBench evaluates
the \emph{finding-text-to-report} task --- given an exam descriptor and concise
findings, produce a faithful report --- and does \emph{not} evaluate primary
image interpretation from DICOM. It therefore measures report faithfulness,
transcription safety, and run-to-run reliability, not whether a model can read a
scan. LAIBench combines (i) a five-dimensional weighted score with a
\emph{deterministic critical-finding hard veto} that forces a failing verdict
regardless of average quality; (ii) a planned \passk{} reliability headline that
asks whether a system gets every critical finding right on all $k$ independent
samples, addressing the saturation and copy-everything gaming that a single-pass
veto rate invites; (iii) eight declared adversarial perturbation classes that
act as metric-validation-by-construction (each perturbation ships with the
dimension and severity the scorer must flag); (iv) a reproducibility hash chain
linking case, suite, scoring code, run, and leaderboard; and (v) inter-rater
agreement, paired-bootstrap discrimination, and judge-calibration tooling. We
are equally explicit about what is \emph{not} yet validated: the critical-finding
detector is currently keyword/substring-based, the public suite is a small
synthetic demo, and LLM-judge reliability against radiologists has not been
measured. We position LAIBench against GREEN, RadGraph/RadGraph-F1, CheXbert,
RadCliQ, ReXVal, RaTEScore, RadFact, and ReXrank, stating for each dimension what
LAIBench borrows and what it adds, and we lay out a validation roadmap centered
on a validated finding extractor and a radiologist-correlation study.
\end{abstract}

\paragraph{Keywords.} radiology report generation; benchmark; clinical safety;
governance; critical findings; reliability; LLM-as-a-judge; reproducibility;
adversarial evaluation.

\section{Introduction}

Automated radiology report generation is moving from research demonstrations to
products that draft or transcribe reports inside clinical workflows. The
dominant evaluation paradigm, inherited from machine translation and
summarization, scores a generated report against a reference and reports an
aggregate number: lexical overlap (BLEU, ROUGE \cite{Papineni2002BLEU,
Lin2004ROUGE}), label agreement (CheXbert \cite{Smit2020CheXbert,
Irvin2019CheXpert}), entity-and-relation overlap (RadGraph/RadGraph-F1
\cite{Jain2021RadGraph}), a learned composite (RadCliQ \cite{Yu2023RadCliQ}), or
a model-graded clinical error count (GREEN \cite{Ostmeier2024GREEN}, RaTEScore
\cite{Zhao2024RaTEScore}). These metrics have advanced the field and correlate,
to varying degrees, with radiologist error judgments \cite{Yu2023ReXVal}.

They leave a governance gap. The question a hospital safety committee or a
regulator asks before turning a system on is not ``what is the mean quality
score across the test set?'' It is closer to: ``If I hand this system the
findings my radiologist dictated, will it ever drop the line that says
\emph{acute intracranial hemorrhage}? Will it ever silently delete a negation
and turn \emph{no fracture} into \emph{fracture}? And will it do the right thing
\emph{every} time, not just on average?'' A high mean score is compatible with a
small but non-zero rate of exactly these catastrophic, individually
unacceptable failures, because averaging dilutes rare critical errors into a
sea of correctly transcribed routine sentences.

LAIBench is built around that governance question. Its design commitments are:

\begin{enumerate}[leftmargin=*]
\item \textbf{Faithfulness over fluency.} The unit of evaluation is whether the
report faithfully preserves the handed-in evidence, not whether it reads well.
\item \textbf{A hard veto, not a soft penalty.} A dropped or inverted critical
finding forces an overall \dimname{FAIL}, independent of how high the weighted
average is. This is a categorical gate, not a deduction.
\item \textbf{Reliability as the headline.} We move from ``did it pass once?'' to
``did it pass $k$ out of $k$ times?'' via a \passk{} metric, because a
single-pass veto rate saturates and is gameable.
\item \textbf{Metric validation by construction.} Eight adversarial perturbation
classes inject known failures with declared expected dimension and severity; the
scorer must catch them, which tests the metric rather than the model.
\item \textbf{Reproducibility and integrity.} A cryptographic hash chain pins
cases, scoring code, runs, and leaderboard, and anti-gaming logic recomputes
every published statistic from raw artifacts.
\end{enumerate}

We stress at the outset that LAIBench is a \emph{benchmark framework}, not a
medical device, regulatory clearance, or clinical validation, and that the
current public artifacts run only on a small synthetic demo suite. The
contribution is a measurement methodology and a set of governance-oriented
instruments, presented so that a skeptical reviewer can audit exactly what is
and is not substantiated.

\section{Scope and Task Definition}
\label{sec:scope}

\paragraph{Task.} LAIBench evaluates the \emph{finding-text-to-report} task.
Each case provides a structured input
\[
x = (\textit{exam descriptor},\ \textit{concise findings},\ \textit{locale})
\]
and the system must emit a complete radiology-style report $\hat{y}$ that
faithfully renders the supplied evidence. The reference target $y$ is a gold
report (or gold structured-finding labels) for the same input. Formally the
benchmark scores a function of $(x, \hat{y}, y)$; it does not consume pixels.

\paragraph{Explicit exclusion.} \textbf{LAIBench does not evaluate primary image
interpretation.} It is given the findings as text and asks whether the generated
report preserves them. It therefore cannot, and does not, measure whether a model
correctly \emph{detected} a finding on a DICOM study. A finding that a human
never dictated cannot be ``missed'' by LAIBench, because LAIBench never sees the
image. This is a deliberate scope boundary, stated up front rather than buried:
LAIBench measures \emph{report faithfulness} and \emph{transcription safety},
not diagnostic image accuracy.

\paragraph{Why this framing, and what it is \emph{not}.} Because the failure mode
LAIBench can detect is a dropped, inverted, or fabricated \emph{handed-in}
finding --- not a finding missed on the image --- we deliberately avoid
``malpractice'' or ``patient-harm'' language. The gate catches a report that
fails to carry forward evidence it was given; whether that constitutes harm
depends on the surrounding clinical workflow, which LAIBench does not model. The
honest claim is narrower and more defensible: the benchmark quantifies how
faithfully and how reliably a generator transcribes and structures provided
clinical evidence.

\paragraph{Output contract.} The public harness expects a constrained
HTML-structured report (titles, sections, line breaks) so that structural
compliance and section ordering are machine-checkable. The contract is a
property of the harness, not a clinical requirement; the scoring dimensions that
matter for safety (critical-finding preservation, negation integrity) are
defined on the report text after tag stripping and normalization.

\paragraph{Language coverage.} The framework leads in English for positioning and
comparability with prior art, while treating Portuguese (pt-BR) as a
\emph{coverage axis} rather than the headline identity. To our knowledge no
public pt-BR report-generation benchmark exists, so the locale dimension extends
coverage into an under-served language; it is not the paper's central claim. The
only public cases shipped today are a small synthetic pt-BR demo suite.

\section{Related Work and Positioning}
\label{sec:related}

Radiology report evaluation has progressed from generic text-overlap metrics
toward clinically grounded ones. \textbf{CheXbert} \cite{Smit2020CheXbert}
labels reports against the CheXpert observation schema \cite{Irvin2019CheXpert}
and compares predicted vs. reference labels, capturing presence/absence of a
fixed set of observations. \textbf{RadGraph} \cite{Jain2021RadGraph} extracts
clinical entities and relations; \textbf{RadGraph-F1} scores entity/relation
overlap against the reference graph. \textbf{RadCliQ} \cite{Yu2023RadCliQ}
combines RadGraph-F1 and BLEU into a single composite calibrated against
radiologist error counts, with the \textbf{ReXVal} \cite{Yu2023ReXVal}
error-annotation dataset as its human ground truth. \textbf{RaTEScore}
\cite{Zhao2024RaTEScore} is an entity-aware similarity metric robust to synonymy
and negation. \textbf{GREEN} \cite{Ostmeier2024GREEN} uses a language model to
enumerate and categorize clinically significant errors, producing an
interpretable error notation and a soft severity-weighted score.
\textbf{RadFact} \cite{Bannur2024RadFact} verifies grounded, sentence-level
factual entailment between generated and reference findings.
\textbf{ReXrank} \cite{Zhang2024ReXrank} is a public leaderboard standardizing
chest X-ray report-generation comparison.

LAIBench is not a replacement for these metrics; it borrows heavily from them
and recombines them under a governance objective. Table~\ref{tab:related} maps
each LAIBench dimension/feature to its closest prior art and states what is
borrowed versus what is added. The central additions are: (1) a \emph{hard veto}
on critical findings rather than a soft penalty; (2) a \passk{} \emph{reliability}
axis layered over per-pass scoring; (3) \emph{adversarial perturbation
validation} of the scorer itself; and (4) a \emph{reproducibility/anti-gaming
governance layer}. The central honest caveat is that several LAIBench dimensions
currently approximate, with keyword and substring heuristics, what the prior-art
metrics implement with trained extractors --- validation of those approximations
is pending (Section~\ref{sec:limitations}).

\begin{table}[t]
\centering
\small
\setlength{\tabcolsep}{4pt}
\renewcommand{\arraystretch}{1.25}
\begin{tabular}{@{}p{2.2cm}p{4.3cm}p{6.0cm}@{}}
\toprule
\textbf{LAIBench feature} & \textbf{Closest prior art} & \textbf{Borrows vs.\ adds} \\
\midrule
\dimname{CRIT} (critical-finding hard veto) &
GREEN soft severity scoring \cite{Ostmeier2024GREEN}; CheXbert label match
\cite{Smit2020CheXbert} &
\emph{Borrows} the idea of severity-graded clinical errors. \emph{Adds} a
categorical hard veto: any dropped/inverted critical finding forces overall
\dimname{FAIL}, which is absent from GREEN's continuous severity score. \\
\dimname{QUAL} (finding preservation, hallucination) &
RadGraph-F1 entity overlap \cite{Jain2021RadGraph}; RadFact entailment
\cite{Bannur2024RadFact}; RaTEScore \cite{Zhao2024RaTEScore} &
\emph{Borrows} entity-recall and hallucination framing. \emph{Approximates}
entity recall via synonym-expanded keyword overlap rather than a trained
extractor --- validation pending. \\
\dimname{TERM} (terminology/locale) &
RadLex-style lexicons; RaTEScore synonym handling \cite{Zhao2024RaTEScore} &
\emph{Adds} locale- and modality-specific terminology checks (incl.\ pt-BR);
rule-based, not a learned model. \\
\dimname{GUIDE} (guideline coverage) &
Structured-reporting / CDE standards \cite{ESR2023Structured} &
\emph{Adds} guideline-system modules (BI-RADS, TI-RADS, LI-RADS, PI-RADS,
Bosniak, Lung-RADS, Fleischner) checking applicability/presence/correctness; not
present in overlap metrics. \\
\dimname{RAG} (evidence fidelity / retrieval) &
IR metrics: P@k, R@k, MRR, nDCG \cite{JarvelinKekalainen2002} &
\emph{Borrows} standard IR metrics for the optional retrieval pipeline.
\emph{Adds} evidence-fidelity structural checks (laterality, levels,
measurements, order) when no retrieval gold exists. \\
Aggregate composite &
RadCliQ learned composite \cite{Yu2023RadCliQ} &
\emph{Differs}: fixed, documented weights plus a hard gate, rather than a
regression fit, for auditability. \\
\passk{} reliability &
$\mathrm{pass@}k$ in code generation \cite{Chen2021Codex} &
\emph{Borrows} the multi-sample reliability idea from program synthesis.
\emph{Adapts} it to ``every critical finding correct on all $k$ samples.'' \\
Adversarial validation &
--- (no direct analogue in report metrics) &
\emph{Adds} eight declared perturbation classes as metric-validation-by-
construction. \\
Leaderboard / provenance &
ReXrank leaderboard \cite{Zhang2024ReXrank} &
\emph{Adds} a cryptographic hash chain and recompute-from-raw anti-gaming
checks. \\
\bottomrule
\end{tabular}
\caption{Positioning of each LAIBench dimension/feature against the closest
prior radiology report-evaluation work. ``Borrows'' marks reused ideas;
``adds''/``approximates'' marks LAIBench-specific design, including approximations
whose validation is pending.}
\label{tab:related}
\end{table}

\section{Scoring Model}
\label{sec:scoring}

\subsection{Five dimensions and weights}

A run scores each case along five dimensions on a $0$--$100$ scale. The default
(research/leaderboard) weights are fixed and documented:
\[
w_{\dimname{CRIT}}=0.30,\quad
w_{\dimname{QUAL}}=0.25,\quad
w_{\dimname{TERM}}=0.20,\quad
w_{\dimname{GUIDE}}=0.15,\quad
w_{\dimname{RAG}}=0.10 .
\]
The dimensions measure:

\begin{description}[leftmargin=1.2cm,style=nextline]
\item[\dimname{CRIT} (0.30)] Critical-finding preservation and unsafe-negation
integrity. With gold critical-finding labels the evaluator computes
recall/precision/F1 against the report, scoring recall-weighted
($\text{score}=70\cdot\text{recall}+30\cdot\text{precision}$) because a missed
critical finding is worse than a false positive; each missed label is emitted as
a separate \texttt{critical}-severity check. A substring/token match inside a
\emph{negated} sentence is counted as a miss, not a hit. Without gold labels the
dimension falls back to structural banned-phrase checks.
\item[\dimname{QUAL} (0.25)] Clinical-finding preservation and hallucination
resistance. Gold findings are matched as exact/partial/missed with a
30-group radiology synonym table (pt-BR + English), and extracted report findings
unmatched to any gold finding are flagged as potential hallucinations (with
pertinent negatives explicitly excluded from the hallucination set). A clinical-
utility floor penalizes missed critical findings.
\item[\dimname{TERM} (0.20)] Locale- and modality-specific terminology and
section/report conventions (e.g.\ correct ultrasound/MRI vocabulary, no
contrast language in a non-contrast exam). Implemented as deterministic
structural rules.
\item[\dimname{GUIDE} (0.15)] Guideline and anatomical coverage via a registry of
modules (Fleischner, BI-RADS, TI-RADS, LI-RADS, PI-RADS, Bosniak, Lung-RADS),
each testing applicability $\rightarrow$ presence $\rightarrow$ correctness, with
valid-value-set validation per system; falls back to anatomy coverage when no
guideline applies.
\item[\dimname{RAG} (0.10)] Evidence fidelity. With retrieval gold it computes
Precision@$k$, Recall@$k$, MRR, and nDCG \cite{JarvelinKekalainen2002} over the
retrieved set (recall-weighted composite); otherwise it scores structural
evidence checks (title tokens, laterality, levels, measurements, section order),
and is marked \dimname{UNSCORED} when neither applies.
\end{description}

When a dimension has no applicable checks it is \dimname{UNSCORED} and excluded;
the remaining weights are renormalized to sum to one, so the composite is always
a proper weighted mean over scored dimensions.

\subsection{Per-dimension severity and the deterministic gate}

Within a dimension, the deterministic score starts as the pass fraction of its
checks. Severity then caps the score: a single critical-severity failure caps the
dimension near the failing range (and additional critical failures lower the cap
further), while three or more major failures cap it at $70$. The dimension
\emph{verdict} is \dimname{FAIL} whenever a critical check fails, regardless of
how many other checks pass --- this is the per-dimension expression of the hard
veto.

\subsection{Combining deterministic and judge scores}

When the optional adversarial LLM judge is enabled, each dimension has a
deterministic score and a judge score (the judge emits $0$--$100$, or $1$--$5$
rescaled to $0$--$100$). The combination is governed by the scoring mode:

\begin{itemize}[leftmargin=*]
\item \textbf{conservative-min} (default for regression calibration): the
combined per-dimension score is $\min(\text{det},\text{judge})$. This is
deliberately pessimistic --- a dimension is only as good as the lower of the two
independent assessments, so the judge cannot rescue a deterministic failure and
vice versa.
\item \textbf{judge-primary}: the judge's $0$--$100$ score is the primary
report-quality value while deterministic critical checks remain gates.
\end{itemize}

The composite overall is the weight-renormalized mean of the combined
per-dimension scores. The overall \emph{verdict} is computed as:
\[
\text{verdict}=
\begin{cases}
\dimname{FAIL} & \text{if any deterministic \emph{or} judge critical failure (when } \texttt{criticalFailForces}),\\
\dimname{PASS} & \text{else if } \text{overall}\ge \tau_{\text{pass}},\\
\dimname{PARTIAL} & \text{else if } \text{overall}\ge \tau_{\text{partial}},\\
\dimname{FAIL} & \text{otherwise.}
\end{cases}
\]
The hard veto is the first clause: a critical failure forces \dimname{FAIL}
\emph{before} the composite is even consulted. This is the property that
distinguishes LAIBench from a soft severity-weighted score --- no quantity of
correctly transcribed routine sentences can lift a report that dropped a
critical finding out of \dimname{FAIL}. The combiner also returns explicit
\texttt{gateReasons} (e.g.\ ``deterministic critical failure'', ``adversarial
critical failure'', ``adversarial phase unavailable'') and a confidence label
that degrades to \texttt{low} when the judge phase is unavailable.

\subsection{Policy profiles}

Three policy profiles parameterize weights and thresholds
(Table~\ref{tab:policies}). \textbf{research} encodes the canonical default
weights with $\tau_{\text{pass}}=84$, $\tau_{\text{partial}}=60$.
\textbf{leaderboard} uses the same canonical weights paired with stricter
submission validation. \textbf{strict} up-weights \dimname{CRIT} to $0.35$ and
raises $\tau_{\text{pass}}=90$, $\tau_{\text{partial}}=70$ for clinical-validation
contexts where critical findings should dominate. In all three,
\texttt{criticalFailForces} is \texttt{true}, so the hard veto is never disabled;
profiles change only how the \emph{non-vetoed} composite is computed and where
the PASS/PARTIAL cutoffs sit. The default (no policy) path and the
\textbf{research} profile produce identical numbers by construction.

\begin{table}[t]
\centering
\small
\begin{tabular}{@{}lccccccc@{}}
\toprule
Profile & \dimname{CRIT} & \dimname{QUAL} & \dimname{TERM} & \dimname{GUIDE} &
\dimname{RAG} & $\tau_{\text{pass}}$ & $\tau_{\text{partial}}$ \\
\midrule
research / leaderboard & 0.30 & 0.25 & 0.20 & 0.15 & 0.10 & 84 & 60 \\
strict & 0.35 & 0.25 & 0.15 & 0.15 & 0.10 & 90 & 70 \\
\bottomrule
\end{tabular}
\caption{Policy profiles. Weights are normalized to sum to one. All profiles keep
the critical-finding hard veto enabled.}
\label{tab:policies}
\end{table}

\section{\texorpdfstring{$\mathrm{pass}^{k}$}{pass-k} Reliability}
\label{sec:passk}

\paragraph{Why the veto rate alone is insufficient.} The deterministic critical
gate yields, per system, a \emph{veto pass-rate}: the fraction of cases with no
critical-finding failure. As a single number this metric \emph{saturates}: a
cautious system that copies essentially all handed-in text verbatim will rarely
drop a critical finding and can score a near-perfect veto pass-rate while being
clinically useless (verbose, unstructured, hallucination-prone elsewhere). The
veto rate is therefore necessary but gameable, and it does not capture the
\emph{run-to-run} stochasticity of sampling-based generators: a system that
preserves a critical finding $90\%$ of the time will pass a single-pass
evaluation $90\%$ of the time but is unacceptable for deployment.

\paragraph{The metric.} Borrowing from $\mathrm{pass@}k$ in program synthesis
\cite{Chen2021Codex}, we define a \emph{reliability} headline. Sample $k$
independent reports per case from a system at its production temperature. Let a
case be \emph{$k$-reliable} if all $k$ samples clear the critical-finding gate
(no dropped, inverted, or fabricated critical finding). Then
\[
\passk \;=\; \frac{1}{|C|}\sum_{c\in C}\mathbf{1}\!\left[\textstyle\bigwedge_{i=1}^{k}\,\text{gate}(\hat{y}_{c,i})=\text{pass}\right],
\]
the fraction of cases for which the system ``gets every critical finding right
$k$ out of $k$ times.'' Unlike the mean veto rate, \passk{} is monotonically
non-increasing in $k$ and exposes systems whose single-pass success hides
unreliable sampling: a per-case gate probability of $0.9$ yields a
$k{=}5$ reliability of $\approx0.59$ under independence, a sharp and honest
reliability signal. Reporting \passk{} alongside the composite score keeps the
deterministic veto as the underlying \emph{gate} while making reliability, not
average quality, the headline.

\paragraph{The critical-safe unit is intentionally narrow.} The per-(case, run)
unit that \passk{} aggregates is \emph{critical-finding safety}, defined more
narrowly than the overall verdict: a run is critical-safe for a case iff it has
\emph{no} deterministic \dimname{CRIT} critical-finding failure \emph{and} no
judge-flagged critical failure. A case can therefore \dimname{FAIL} its overall
verdict on weak \dimname{RAG} or structure while still being critical-safe; that
distinction is the entire point of a critical-finding safety metric, so \passk{}
isolates the safety-relevant gate from general quality. The harness also reports a
companion \passk{} over the overall verdict (verdict $\neq$ \dimname{FAIL} on all
$k$ runs) for completeness.

\paragraph{Implementation status.} The \passk{} meta-metric is now implemented in
the harness (\texttt{src/reliability.ts}): it reuses the existing per-sample
critical gate, defines the critical-safe unit above, and aggregates over $k$
independent runs of the same system on the same suite without altering scoring.
It is not yet wired as the default leaderboard \emph{headline}, and --- as with
every result in this paper --- we report \emph{no} measured \passk{} values,
because there are no real-data runs (Section~\ref{sec:limitations}). We describe
\passk{} as the intended reliability headline, computed by available code but not
yet validated on clinical data.

\section{Adversarial Perturbation Validation}
\label{sec:perturb}

A benchmark's credibility rests on whether its scorer actually catches the
failures it claims to penalize. LAIBench tests this \emph{by construction}: it
ships eight perturbation generators, each of which takes a (case, gold report)
pair and produces a deliberately defective output that the scorer
\emph{must} flag. Crucially, every perturbation declares, in code, the
dimension(s) and severity it is supposed to trigger (Table~\ref{tab:perturb}).
This turns metric validation into a falsifiable test: if a perturbation that
declares a \texttt{critical} \dimname{CRIT} failure is not caught, the metric ---
not the model --- has a bug.

\begin{table}[t]
\centering
\small
\setlength{\tabcolsep}{5pt}
\renewcommand{\arraystretch}{1.2}
\begin{tabular}{@{}llll@{}}
\toprule
\textbf{Perturbation} & \textbf{Expected dim(s)} & \textbf{Severity} & \textbf{What it injects} \\
\midrule
\texttt{laterality\_flip}      & \dimname{RAG}, \dimname{CRIT} & critical & swaps left/right (esquerdo/direito) \\
\texttt{negation\_drop}        & \dimname{CRIT}                & critical & removes a negation (``no fracture''$\rightarrow$``fracture'') \\
\texttt{negation\_insert}      & \dimname{CRIT}                & critical & inserts a false negation before a present finding \\
\texttt{measurement\_scramble} & \dimname{RAG}, \dimname{QUAL} & major    & replaces measurements with implausible values \\
\texttt{critical\_drop}        & \dimname{CRIT}                & critical & deletes every declared critical finding \\
\texttt{critical\_invent}      & \dimname{CRIT}                & critical & adds a fabricated critical finding \\
\texttt{terminology\_corrupt}  & \dimname{TERM}, \dimname{QUAL}& major    & swaps canonical terms for colloquialisms ($\ge$30 rules/locale) \\
\texttt{structure\_break}      & \dimname{QUAL}, \dimname{RAG} & major    & strips required HTML structure and section labels \\
\bottomrule
\end{tabular}
\caption{The eight adversarial perturbation classes with their \emph{declared}
expected dimension and severity. Each is deterministic given (caseId, kind) via a
seeded PRNG, so the validation suite is reproducible.}
\label{tab:perturb}
\end{table}

\paragraph{Catch rule.} A perturbation is \emph{caught} if, for any of its
expected dimensions $D$ at expected severity $S$, the scored result shows
\emph{either} (a) a deterministic check on $D$ failing at severity $\ge S$,
\emph{or} (b) a judge critical failure on $D$, \emph{or} (c) a combined
$D$-dimension score below a per-severity floor (critical $<60$, major $<80$,
minor $<90$). The harness aggregates per-kind and overall \emph{catch rates} and
labels the scorer \emph{robust} ($\ge90\%$), \emph{leaky} ($\ge70\%$), or
\emph{broken} otherwise. Because the perturbations are deterministic and the
expected dimensions are declared, the catch-rate report is a unit test of the
metric: a regression that weakens the negation handler will show up as a drop in
the \texttt{negation\_drop}/\texttt{negation\_insert} catch rate, not as a silent
loss of clinical safety.

\paragraph{Scope of this validation.} This is validation of the \emph{scorer's
sensitivity to known defects}, not validation against radiologist judgment. It
demonstrates that the metric reacts as designed to constructed failures; it does
not demonstrate that the metric's notion of a ``critical finding'' matches a
radiologist's. The latter is the subject of Section~\ref{sec:limitations}'s
roadmap.

\section{Statistical Methodology}
\label{sec:stats}

LAIBench ships the statistical instruments needed to treat benchmark outputs as
measurements with uncertainty rather than point estimates.

\paragraph{Inter-rater agreement.} For gold-label construction and judge
consistency we implement Cohen's $\kappa$ \cite{Cohen1960Kappa} (two raters,
nominal), Fleiss' $\kappa$ \cite{Fleiss1971Kappa} ($N$ raters, nominal), and
Krippendorff's $\alpha$ \cite{HayesKrippendorff2007,Krippendorff2018} for
interval data ($0$--$100$ scores, with missing-value support). We report
Landis--Koch \cite{LandisKoch1977} interpretive bands for $\kappa$. The interval
$\alpha$ implementation was recently corrected to the canonical ordered-pair
formula: observed disagreement counts each within-unit value pair in both orders
($i\neq j$), consistent with the expected-disagreement normalization
$2/(N(N-1))$ over the value pool. The earlier unordered-pair counting halved
$D_o$ and inflated $\alpha$ via $\alpha_{\text{wrong}}=(\alpha_{\text{true}}+1)/2$;
the corrected routine is checked against the R \texttt{irr}/\texttt{krippendorff}
implementations.

\paragraph{Discrimination via paired bootstrap.} To test whether the benchmark
separates two systems on the same locked suite, we compute per-case score
differences and a seeded paired bootstrap \cite{EfronTibshirani1993} (default
$10{,}000$ resamples) yielding the mean difference, a percentile confidence
interval, and a two-sided $p$-value. A separation is declared significant only
under a \emph{double gate}: $p<\alpha$ \emph{and} the CI excludes zero (lower
bound $>0$). The harness further requires an effect-size floor (default
$\ge5$ percentage points) to call a separation ``meaningful'' versus merely
``statistically significant,'' and it stratifies by modality and difficulty
(strata with $n\ge5$) to flag any stratum that reverses the global ranking. A
reference probe additionally scores each gold report \emph{as if it were a model
output}; gold should score near-perfect, so low gold scores expose
dataset/scoring bugs rather than model weakness.

\paragraph{Judge calibration.} Because the LLM judge is optional but
consequential, we provide a calibration suite computing: (i) \emph{test-retest}
reliability (same judge, repeated runs on the same suite) as Krippendorff
$\alpha$ over continuous overall scores; (ii) \emph{cross-judge} agreement
(different judge models) as Cohen's $\kappa$ on verdicts and $\alpha$ on overall
scores; and (iii) a \emph{det-vs-judge} Spearman correlation
\cite{Spearman1904} testing whether judge scores track deterministic scores in
the expected direction. A judge is labeled \emph{calibrated}/\emph{weak}/
\emph{uncalibrated} from these signals. This is standard LLM-as-a-judge hygiene
\cite{Zheng2023JudgeMTBench}.

\paragraph{Honest gap.} These instruments measure judge \emph{self-consistency}
and \emph{internal} det-vs-judge coherence. They do \emph{not} measure judge
reliability against radiologists, which we have not yet done. A judge that is
perfectly self-consistent and perfectly correlated with the deterministic scorer
could still disagree systematically with expert reviewers; closing that gap
requires the radiologist study in Section~\ref{sec:limitations}.

\section{Reproducibility and Provenance}
\label{sec:provenance}

Every comparable result is backed by a cryptographic hash chain so that anyone
can verify a leaderboard row was produced from the exact cases, scoring code, and
judge configuration claimed:
\[
\text{caseHash}\ \rightarrow\ \text{suiteHash}\ \rightarrow\ \text{scoringHash}\ \rightarrow\ \text{runHash}\ \rightarrow\ \text{leaderboardHash}.
\]
A \textbf{caseHash} is the SHA-256 of the canonicalized case (id, exam, findings,
locale, gold findings, critical findings, guideline expectations, context,
difficulty). A \textbf{suiteHash} is the SHA-256 of the sorted case hashes. A
\textbf{scoringHash} is the SHA-256 over a pinned, enumerated list of scoring and
evaluator source files; building the provenance manifest \emph{fails loudly} if
any listed file is missing, so a stale file list cannot silently invalidate the
chain. A \textbf{runHash} binds the suite hash, the canonicalized run manifest
(benchmark version, suite, locale, track, provider, model label, scaffold, judge
provider/model, submission mode, canary token), and the scoring hash. A
\textbf{leaderboardHash} is the SHA-256 of the sorted run hashes. Canonical JSON
serialization (sorted keys, deterministic array encoding) makes every hash
order-insensitive and reproducible. Perturbations and the bootstrap use seeded
PRNGs, so the validation and discrimination outputs are bit-reproducible across
machines.

\section{Leaderboard Integrity and Anti-Gaming}
\label{sec:integrity}

A public leaderboard is only as trustworthy as its resistance to inflated
self-reports. LAIBench enforces several integrity properties:

\begin{itemize}[leftmargin=*]
\item \textbf{Recompute, never trust.} Leaderboard and compare operations
recompute case overalls, suite summaries, verdict counts, per-dimension means,
comparable keys, and suite hashes from raw artifacts. If a run JSON was hand-
edited to inflate \texttt{averageOverall}, strict-pass rate, non-fail rate,
per-dimension means, or comparable metadata, leaderboard generation aborts before
producing public output.
\item \textbf{Comparable keys.} Paired comparisons are only permitted between
runs sharing a comparable key (same suite hash, locale, track); the framework
refuses to rank incomparable runs as equivalent and refuses superiority claims
without paired testing on the same suite hash.
\item \textbf{Track separation.} Product agents, custom agents, mini-agent
scaffolds, and raw models occupy separate ranks rather than being merged.
\item \textbf{Contamination scanning.} A per-suite canary token is searched
whitespace-insensitively across raw, normalized, and sanitized report text
(defeating split-token evasion); the judge can also flag memorized boilerplate.
Runs are labeled \emph{clean}/\emph{suspicious}/\emph{contaminated}.
\item \textbf{Copy-everything is not a win.} Because \dimname{TERM},
\dimname{GUIDE}, structural compliance, and the hallucination component of
\dimname{QUAL} all penalize verbose, unstructured, or unsupported content, a
report that copies all handed-in text to maximize the critical-finding veto rate
is penalized elsewhere --- and \passk{} (Section~\ref{sec:passk}) further removes
the incentive to game a single-pass veto.
\end{itemize}

\section{Threats to Validity and Limitations}
\label{sec:limitations}

We list the threats that a reviewer should weigh; several are foundational and
bound the claims this paper can make.

\begin{description}[leftmargin=0.6cm,style=nextline]
\item[Text, not image.] LAIBench evaluates report faithfulness from provided
text evidence and does \emph{not} evaluate primary image interpretation
(Section~\ref{sec:scope}). It cannot detect findings missed on the scan and
makes no diagnostic-accuracy claim.
\item[Keyword/substring critical detector.] The critical-finding detector is
currently keyword/substring- and token-overlap-based with negation handling. It
can miss phrasing variants and is not a validated clinical extractor. The
intended remedy is a validated extractor (e.g.\ a GREEN-based
\cite{Ostmeier2024GREEN} or RadGraph-based \cite{Jain2021RadGraph} entity model)
in place of the heuristic. Until then, \dimname{CRIT}/\dimname{QUAL} should be
read as \emph{approximations} of recall, not validated clinical recall.
\item[Heuristic gold and approximate entity recall.] \dimname{QUAL} approximates
entity recall via synonym-expanded keyword overlap rather than a trained
extractor, and gold findings in the daily suites are heuristically derived from
reports rather than fully radiologist-adjudicated. Both approximations are
documented and both are pending validation.
\item[Synthetic-only public data; no real-data results yet.] The only public
cases are a small synthetic pt-BR demo suite for smoke tests and framework
review. \textbf{This paper reports no results on real clinical data.} All current
numbers (catch rates, discrimination, calibration) are illustrative of the
machinery on the synthetic demo suite; they are not evidence about any deployed
system's clinical performance.
\item[Unmeasured judge-vs-radiologist reliability.] We measure judge self-
consistency and det-vs-judge coherence, not agreement with radiologists
(Section~\ref{sec:stats}). The LLM-as-a-judge component is optional and its
clinical validity is unestablished.
\item[Vendor-built.] LAIBench is developed by a commercial radiology-reporting
vendor. We mitigate the conflict by publishing the scoring code, fixing weights
rather than fitting them, providing recompute-from-raw anti-gaming checks, and
documenting claim discipline; but the benchmark has not been independently
governed, and a system built by the same vendor should not be ranked on it
without third-party oversight.
\item[Constrained output contract.] The HTML structural contract aids machine
checking but is not a clinical standard; structural-compliance dimensions
(\dimname{TERM}, parts of \dimname{RAG}) reflect that contract, not necessarily
local reporting requirements.
\end{description}

\paragraph{Validation roadmap.} The path from ``engineering benchmark readiness''
to defensible clinical claims is: (1) replace the keyword critical detector with
a validated extractor and re-run perturbation validation against it; (2) lock a
radiologist-adjudicated subset (the framework already specifies a two-reviewer,
$\ge80\%$ acceptability / $\ge70\%$ per-dimension agreement adjudication gate);
(3) measure judge-vs-radiologist agreement and calibrate the judge to that ground
truth; (4) run paired discrimination on real, held-out clinical suites that are
never used for system iteration; and (5) only then upgrade public wording beyond
engineering readiness. Until these are done, LAIBench claims report faithfulness,
transcription safety, and reliability measurement --- not prospective clinical
safety.

\section{Conclusion}

LAIBench reframes radiology report-generation evaluation around the question a
deployer actually asks: not ``how good is the average report?'' but ``does the
system faithfully preserve every critical finding it was handed, every time?''
The instruments that follow from this reframing --- a critical-finding hard veto
instead of a soft penalty, a \passk{} reliability headline over a saturating veto
rate, eight declared adversarial perturbations that validate the scorer by
construction, a paired-bootstrap discrimination test with a double significance
gate, a judge-calibration suite, and a cryptographic provenance and anti-gaming
layer --- are, taken together, a governance-oriented measurement methodology
rather than another quality metric. We have been deliberate about scope (text,
not image) and about what remains unvalidated (keyword heuristics, synthetic-only
public data, no real-data results, unmeasured judge-vs-radiologist reliability,
vendor authorship). The contribution we claim is the methodology and its honest
boundary conditions; the validation roadmap above is the work that would let the
benchmark support stronger claims.

\section*{Reproducibility}
The scoring code, evaluators, perturbation generators, statistics, provenance
hash chain, schemas, and the synthetic demo suite are available in the public
LAIBench repository. The full clinical corpus, hidden test set, answer keys, and
private scoring criteria are gated and are not part of the public release.

\section*{Conflicts of Interest}
LAIBench is developed by a commercial radiology-reporting vendor. See
Section~\ref{sec:limitations}.

\bibliographystyle{plain}
\bibliography{references}

\end{document}
